\documentclass[12pt,a4paper]{report}
\usepackage[]{amsmath,amsthm,amssymb,amscd}
\usepackage[latin1]{inputenc}
%\usepackage{fullpage}

\topmargin 0pt
\advance \topmargin by -\headheight
\advance \topmargin by -\headsep
     
\textheight 246mm
     
\oddsidemargin 0pt
\evensidemargin \oddsidemargin
\marginparwidth 0.5in
     
\textwidth 159mm


%               Theorem Environments
\theoremstyle{definition}
\newtheorem{ex}{}
\newcommand{\pd}[1]{\frac{\partial}{\partial #1}} %\pd{x}
%               Subquestions numbered with letters
\renewcommand{\theenumi}{\textbf{\alph{enumi}}}

%               Maths definitions

\newcommand{\NN}{\ensuremath{\mathbb N}}
\newcommand{\ZZ}{\ensuremath{\mathbb Z}}
\newcommand{\CC}{\ensuremath{\mathbb C}}
\newcommand{\TT}{\ensuremath{\mathbb T}}
\newcommand{\RR}{\ensuremath{\mathbb R}}
\newcommand{\QQ}{\ensuremath{\mathbb Q}}
\newcommand{\g}{\ensuremath{\frak{g}}}
\newcommand{\h}{\ensuremath{\frak{h}}}
\newcommand{\ft}{\ensuremath{\frak{t}}}
\newcommand{\cB}{\mathcal{B}}
\newcommand{\cN}{\mathcal{N}}
\newcommand{\cL}{\mathcal{L}}
\newcommand{\cD}{\mathcal{D}}
%\newcommand{\pd}[1]{\frac{\partial}{\partial #1}} %\pd{x}
%\newcommand{\pd}[1]{\frac{\partial}{\partial #1}} %\pd{x}

\begin{document}
\noindent
{\sffamily\large Marco Zambon 
\hfill     29/03/2011\\Geometr\'ia III, curso 2010-11}
 

%\noindent
%\hrulefill{}\\
\begin{center}\bfseries\large\textsf{Examen Parcial    \vspace{-2mm}}
\end{center}

 
%\hrulefill{}\\

 \noindent
\hrulefill{}\\
\textsf{Por favor, escribid en cada hoja vuestro nombre y el n\'umero del problema.\\
Pod\'eis utilizar, cit\'andolos de manera apropriada, resultados que hemos obtenido en clase o en los ejercicios para casa.    \vspace{-2mm}}

\noindent
\hrulefill{}\\

 
\begin{ex}[8 puntos]    
Considera la variedad topol\'ogica con borde $$M:=\{(x,y):x,y \in \RR,  1 \le \sqrt{x^2+y^2}\le 2\}$$
y considera esta relaci\'on de equivalencia   en $M$: cada punto de $M$ es equivalente s\'olo a s\'i mismo excepto  los puntos del borde $\{(x,y): \sqrt{x^2+y^2}=1\}\cup \{(x,y): \sqrt{x^2+y^2}=2\}$, para los que se cumple
$$(x,y) \sim (2x,2y)$$
para cada $(x,y) \in M$ con $\sqrt{x^2+y^2}=1$.

Construye un homeomorfismo del cociente  $M/\sim$ a una superficie en
la lista del teorema de clasificaci\'on de superficies compactas y conexas.
\end{ex}

\begin{ex} [7 puntos]
Sea $P_2(\RR)$ el plano proyectivo y $\TT:=S^1\times S^1$ el toro. > A qu\'e superficie en la lista de superficies compactas y conexas es homeomorfa la suma conexa $$P_2(\RR) \; \sharp \;  \TT\;\;?$$  
Demuestra.
\end{ex}
 

\begin{ex}  [8 puntos]
> Es verdad que $\RR$, con su topolog\'ia canonica, admite un numero infinito de estructuras  de variedad diferencial (distintas entre s\'i)? Demu\'estralo.
\end{ex}





\begin{ex}[7 puntos]
Sea $M$ una variedad diferencial de dimensi\'on $n$, $S$  una subvariedad de $M$ de dimensi\'on $k$, y $p$ un punto de $S$. Demuestra que existe un entorno abierto $U$ de $p$ en $M$ y una aplicaci\'on diferenciable $F \colon U \to \RR^{n-k}$ tal que
$$U\cap S= F^{-1}(0).$$
\end{ex}

 

%\begin{ex}  
%Demuestra que 
%$$T \;\;\sharp \;\; C\cong T-(D_1 \cup D_2)$$
%donde $T=S^1\times S^1$ es el toro, $C=S^1\times [0,1]$ el cilindro, y $D_i\subset T$ (por $i=1,2$) son discos cerrados tal que $D_1\cap D_2=\emptyset$.
%\end{ex} 






 \end{document}
 %%%%%%%
 \begin{ex}

\end{ex}

 \begin{ex}
\begin{enumerate}
\item 
\end{enumerate}
\end{ex}
 
 
 
 